\documentclass[11pt]{article}
\usepackage[T1]{fontenc}
\usepackage[utf8]{inputenc}
\usepackage{amsmath,amssymb,amsthm}
\usepackage{hyperref}
\usepackage{mathtools}

\title{Correction (version detaillee)\\Automate residuel et instrumentation}
\author{}
\date{}

\newcommand{\Tr}{\mathcal{T}}
\newcommand{\AP}{\mathsf{AP}}
\newcommand{\eval}{\models}
\newcommand{\Prog}{\mathrm{Prog}}
\newcommand{\nnf}{\mathrm{nnf}}
\newcommand{\res}{\mathrm{res}}
\newcommand{\step}{\mathrm{step}}
\newcommand{\states}{\mathsf{States}}
\newcommand{\deltaa}{\delta}
\newcommand{\atoms}{\mathsf{Atoms}}
\newcommand{\True}{\mathsf{true}}
\newcommand{\False}{\mathsf{false}}

\theoremstyle{definition}
\newtheorem{definition}{Definition}
\newtheorem{lemma}{Lemma}
\newtheorem{theorem}{Theorem}
\newtheorem{corollary}{Corollary}
\newtheorem{proposition}{Proposition}

\begin{document}
\maketitle

\section{Preliminaires}

\subsection{Etats, traces et valuations}

\begin{definition}[Etats et traces]
Soit $\Sigma$ l'ensemble des etats. Une \emph{trace} est une suite infinie
$\tau = s_0 s_1 s_2 \ldots$ avec $s_i \in \Sigma$. On note $\tau[i] = s_i$.
\end{definition}

\begin{definition}[Atomes et valuations]
Soit $\AP$ un ensemble fini d'atomes propositionnels. Une valuation
$\sigma : \AP \to \{\True,\False\}$ associe une valeur a chaque atome.
Chaque etat $s \in \Sigma$ induit une valuation $\sigma_s$.
\end{definition}

\subsection{Fragment LTL $G/X$}

\begin{definition}[Syntaxe]
Les formules sont definies par:
\[
\varphi ::= \True \mid \False \mid p \in \AP \mid \neg\varphi
\mid (\varphi \wedge \varphi) \mid (\varphi \vee \varphi)
\mid X\varphi \mid G\varphi.
\]
On suppose les formules en NNF (negations uniquement sur les atomes).
\paragraph{NNF.}
La transformation $\nnf(\cdot)$ pousse les negations sur les atomes et
rejette les negations au-dessus de $G$ (hors fragment).
\end{definition}

\begin{definition}[Semantique LTL]
Pour une trace $\tau$ et un indice $i$:
\begin{align*}
\tau,i \eval p &\iff \sigma_{\tau[i]}(p) = \True \\
\tau,i \eval \neg p &\iff \sigma_{\tau[i]}(p) = \False \\
\tau,i \eval \varphi \wedge \psi &\iff \tau,i \eval \varphi \text{ et } \tau,i \eval \psi \\
\tau,i \eval \varphi \vee \psi &\iff \tau,i \eval \varphi \text{ ou } \tau,i \eval \psi \\
\tau,i \eval X\varphi &\iff \tau,i+1 \eval \varphi \\
\tau,i \eval G\varphi &\iff \forall j \ge i,\ \tau,j \eval \varphi.
\end{align*}
On note $\tau \eval \varphi$ pour $\tau,0 \eval \varphi$.
\end{definition}

\subsection{Progression}

\begin{definition}[Progression]
La progression $\Prog(\varphi,\sigma)$ renvoie une formule residuelle telle que:
\[
\tau,i \eval \varphi \iff \tau,i+1 \eval \Prog(\varphi,\sigma_{\tau[i]}).
\]
Sur le fragment $G/X$ (NNF):
\[
\Prog(p,\sigma)=\begin{cases}\True & \sigma(p)=\True\\ \False & \text{sinon}\end{cases}
\quad
\Prog(\neg p,\sigma)=\begin{cases}\True & \sigma(p)=\False\\ \False & \text{sinon}\end{cases}
\]
\[
\Prog(\varphi\wedge\psi,\sigma)=\Prog(\varphi,\sigma)\wedge\Prog(\psi,\sigma),\quad
\Prog(\varphi\vee\psi,\sigma)=\Prog(\varphi,\sigma)\vee\Prog(\psi,\sigma)
\]
\[
\Prog(X\varphi,\sigma)=\varphi,\quad
\Prog(G\varphi,\sigma)=\Prog(\varphi,\sigma)\wedge G\varphi.
\]
\end{definition}

\begin{lemma}[Correction de la progression]
Pour toute trace $\tau$, tout $i$ et toute formule $\varphi$:
\[
\tau,i \eval \varphi \iff \tau,i+1 \eval \Prog(\varphi,\sigma_{\tau[i]}).
\]
\end{lemma}
\begin{proof}
Par induction structurelle sur $\varphi$.
\end{proof}

\section{Automate residuel}

\subsection{Residus}

\begin{definition}[Residus]
Soit $\varphi$ en NNF. On definit l'ensemble des residus $\states$ comme
l'ensemble des formules obtenues par applications iteratives de $\Prog$
depuis $\varphi$, puis simplification (idempotence, constantes, flatten).
\end{definition}

\paragraph{Simplification.}
On utilise une simplification booleenne standard:
\[
\True \wedge \psi = \psi,\ \False \wedge \psi = \False,\ 
\True \vee \psi = \True,\ \False \vee \psi = \psi.
\]
Elle preserve la semantique: $\tau,i \eval \mathrm{Simplify}(\varphi) \iff \tau,i \eval \varphi$.

\paragraph{Finitude.}
Dans le fragment $G/X$, la fermeture par sous-formules est finie; la progression
ne produit que des sous-formules (plus $\True/\False$). Ainsi $\states$ est fini.

\begin{definition}[Automate residuel]
L'automate $\mathcal{A}_\varphi$ est:
\[
\mathcal{A}_\varphi = (\states, q_0, \deltaa)
\]
avec $q_0=\varphi$ et
$\deltaa(q,\sigma)=\mathrm{Simplify}(\Prog(q,\sigma))$.
\end{definition}

\begin{lemma}[Invariance du residu]
Soit $\tau$ une trace et $q_i$ la suite d'etats definie par
$q_0=\varphi$, $q_{i+1}=\deltaa(q_i,\sigma_{\tau[i]})$. Alors:
\[
\tau,i \eval q_i \quad \text{pour tout } i.
\]
\end{lemma}
\begin{proof}
Par induction sur $i$ et le lemme de progression.
\end{proof}

\begin{theorem}[Correction de l'automate (version semantique)]
Pour toute trace $\tau$:
\[
\tau \eval \varphi \iff \forall i,\ \tau,i \eval q_i.
\]
\end{theorem}
\begin{proof}
Par le lemme d'invariance du residu.
\end{proof}

\paragraph{Remarque.}
La condition syntactique $q_i \ne \False$ est seulement \emph{sondante} si la
simplification est complete. Dans cette version, on utilise donc la condition
semantique $\tau,i \eval q_i$.

\section{Instrumentation (langage imperatif jouet)}

\subsection{Programme}

\begin{definition}[Systeme de transition]
Un programme est un systeme $(\Sigma,\rightarrow,s_0)$.
Une execution est une trace $\tau$ telle que $s_0 \rightarrow s_1 \rightarrow \cdots$.
\end{definition}

\subsection{Instrumentation moniteur}

\begin{definition}[Etat instrumente]
On etend l'etat a $\Sigma'=\Sigma \times \states$.
Etat initial $(s_0,q_0)$.
Transition:
\[
(s,q)\rightarrow (s',q') \iff s\rightarrow s' \text{ et } q'=\deltaa(q,\sigma_s).
\]
\end{definition}

\begin{definition}[Invariant de securite]
On ajoute l'invariant semantique $\tau,i \eval q_i$ (ou $q$ n'est pas ``mauvais'').
\end{definition}

\begin{theorem}[Correction de l'instrumentation]
Pour toute execution instrumentee $\tau'=(s_0,q_0)(s_1,q_1)\ldots$:
\[
\forall i,\ \tau,i \eval q_i \iff s_0 s_1 s_2 \ldots \eval \varphi.
\]
\end{theorem}
\begin{proof}
La projection $\pi(\tau')$ est une trace du programme. La suite $(q_i)$ suit
exactement l'automate residuel par construction, donc on applique la correction
de l'automate.
\end{proof}

\section{Traduction Why3 (vue logique)}

\subsection{Decoupage des obligations}

\paragraph{Contrats.} Les contrats sont convertis en une formule LTL
$\psi$ pour le moniteur:
\[
\psi = (\bigwedge \text{Assume/Requires}) \Rightarrow (\bigwedge \text{Guarantee/Ensures})
\]
et en obligations locales pour \texttt{step}.

\paragraph{Obligations de \texttt{step}.}
On construit des listes \texttt{pre} et \texttt{post}:
\begin{itemize}
\item \texttt{pre}: obligations issues de Requires/Assume, eventuellement gardees
par \texttt{\_\_first\_step} si la formule n'est pas sous $G$.
\item \texttt{post}: obligations issues de Ensures/Guarantee, gardees par
\texttt{old(\_\_first\_step)} si la formule n'est pas sous $G$, plus les gardes
par transition (\texttt{old(st)=src}) et les invariants/liaisons techniques
(pre, pre\_k, fold, instances, etc.).
\end{itemize}

\subsection{k-guard}
Pour une formule $G\alpha$ avec profondeur $k$ en $X$, on ajoute une garde
(\texttt{\_\_step\_count >= k}) afin de ne requerir l'obligation qu'apres $k$ pas.

\paragraph{Initial-only.}
Si $\varphi$ n'est pas de la forme $G(\cdot)$, alors ses obligations
sont gardees par \texttt{\_\_first\_step} (ou \texttt{old(\_\_first\_step)} en post).

\section{Conclusion}
La correction globale se deduit de:
\begin{enumerate}
\item la correction de l'automate residuel vis-a-vis de $\varphi$;
\item la correction de l'instrumentation (etat etendu + invariant).
\end{enumerate}
Ainsi, la satisfaction de l'invariant semantique implique la satisfaction
de la specification LTL sur toutes les executions.

\end{document}
